%!TEX program = xelatex
\documentclass[10pt, letterpaper]{article}

% ── Packages ──────────────────────────────────────────────────
\usepackage[margin=0.65in, top=0.6in, bottom=0.6in]{geometry}
\usepackage{fontspec}
\setmainfont{Times New Roman}

\usepackage{enumitem}
\usepackage{titlesec}
\usepackage{xcolor}
\usepackage{hyperref}
\usepackage{tabularx}
\usepackage{array}
\usepackage{parskip}


% ── Colors ────────────────────────────────────────────────────
\definecolor{navy}{HTML}{1A3A5C}
\definecolor{linkblue}{HTML}{2E5E9E}
\definecolor{darkgray}{HTML}{333333}
\definecolor{midgray}{HTML}{555555}
\definecolor{lightgray}{HTML}{888888}
\definecolor{rulecol}{HTML}{2E5E9E}

% ── Hyperref ──────────────────────────────────────────────────
\hypersetup{
  colorlinks=true,
  urlcolor=linkblue,
  linkcolor=linkblue,
  pdftitle={Lingzhi Zhao -- CV},
  pdfauthor={Lingzhi Zhao}
}

% ── Fonts ─────────────────────────────────────────────────────
\renewcommand{\familydefault}{\sfdefault}   % sans-serif throughout

% ── Section styling ───────────────────────────────────────────
\titleformat{\section}
  {\large\bfseries\color{navy}}
  {}{0em}{\MakeUppercase}
  [\color{rulecol}\titlerule]
\titlespacing{\section}{0pt}{8pt}{4pt}

% ── List styling ──────────────────────────────────────────────
\setlist[itemize]{
  leftmargin=1.4em,
  itemsep=1pt,
  parsep=0pt,
  topsep=2pt,
  label=\textcolor{navy}{\textbullet}
}

% ── Helpers ───────────────────────────────────────────────────
% Role entry: {Institution}{Role}{Date}
\newcommand{\roleentry}[3]{%
  \vspace{4pt}%
  \noindent
  \begin{tabularx}{\linewidth}{@{}X r@{}}
    \textbf{\color{darkgray}#1} & \textcolor{lightgray}{\small #3}
  \end{tabularx}\\[1pt]
  \noindent\textit{\color{midgray}\small #2}%
  \vspace{2pt}
}

% Project label
\newcommand{\project}[1]{%
  \vspace{2pt}%
  \noindent\textbf{\color{navy}\small #1}%
  \vspace{1pt}
}

% Skills row label
\newcommand{\skills}[1]{\textbf{\color{darkgray}\small #1}}

% Publication entry
\newcommand{\pub}[3]{%
  \vspace{3pt}%
  \noindent\hangindent=1.6em\hangafter=1
  \hypertarget{pub#1}{}%
  \textbf{\color{navy}[#1]}\enspace #2 \textit{\color{midgray}#3}\par
}

% ── No paragraph indent ───────────────────────────────────────
\setlength{\parindent}{0pt}
\setlength{\parskip}{0pt}

%══════════════════════════════════════════════════════════════
\begin{document}
%══════════════════════════════════════════════════════════════

% ── Header ────────────────────────────────────────────────────
\begin{center}
  {\Huge\bfseries\color{navy} Lingzhi Zhao}\\[2pt]
  {\small\color{midgray}
    Ph.D.\ Candidate in Computer Science \,|\, University of Illinois Urbana-Champaign}\\[1pt]
  {\small
    +1 (217) 979-9381 \,|\,
    \href{mailto:lz26@illinois.edu}{lz26@illinois.edu} \,|\,
    \href{https://chutoutian.github.io}{Website} \,|\,
    \href{https://linkedin.com/in/lingzhizhao}{LinkedIn}
  }
\end{center}

%──────────────────────────────────────────────────────────────
\section{Research Interests}
%──────────────────────────────────────────────────────────────
\iffalse
{\small\color{midgray}
  Networked multimedia systems (video streaming, XR, egocentric AI)
  $\cdot$ Intelligent mobile systems $\cdot$ Generative streaming\par}
  \fi

%My research broadly focuses on \textbf{networked multimedia systems (video on demand, live streaming, XR)} and \textbf{intelligent mobile systems}. 
%I am interested in designing efficient systems for emerging media systems, including XR, egocentric AI assistants, AI-mediated communication systems.
%My research lies at the \textbf{intersection of emerging media systems} (e.g., XR) \textbf{and multimodal AI}. I am interested in designing efficient and intelligent systems that bring multimodal understanding to real-world interactive and streaming applications.

My research bridges \textbf{multimodal AI} and \textbf{multimedia systems}. I am interested in adapting vision–language models (VLMs) via post-training and efficient on-device deployment, which brings multimodal understanding to emerging media systems such as XR.

%──────────────────────────────────────────────────────────────
\section{Education}
%──────────────────────────────────────────────────────────────

\roleentry{University of Illinois Urbana-Champaign}
  {Ph.D.\ in Computer Science $\cdot$ Advisor: Prof.\ Klara Nahrstedt}
  {2022 Aug. -- 2026 Dec. (expected) }

\roleentry{Shanghai Jiao Tong University}
  {M.S.\ in Electrical and Computer Engineering $\cdot$ Advisor: Prof.\ Ying Cui}
  {2019 -- 2022}

\roleentry{Shanghai University}
  {B.S.\ in Electrical and Computer Engineering}
  {2015 -- 2019}

%──────────────────────────────────────────────────────────────
\section{Skills}
%──────────────────────────────────────────────────────────────

\begin{tabular}{@{}l@{\hspace{6pt}}l@{\hspace{6pt}}l@{}}
  \skills{Machine Learning}      && PyTorch, Transformers, VLM/MLLM, SFT \& RL Post-training, On-device Inference Optimization \\
  \skills{Tools and Languages}   && Python, C++, MATLAB, Swift, Java, C\#, Rust \\[2pt]
  \skills{(Video) Systems \& Dev.} && FFmpeg, DASH, ABR \& Rate Control, NVCodec, Wireshark, iperf3, tc, Git, NS-3 \\[2pt]

\end{tabular}

%──────────────────────────────────────────────────────────────
\section{Research Experience}
%──────────────────────────────────────────────────────────────

\roleentry{Multimedia Operating Systems and Networking (MONET) Lab, UIUC}
  {Graduate Research Assistant $\cdot$ Advisor: Prof.\ Klara Nahrstedt}
  {Aug.\ 2022 -- Present}

\project{AquaVLM: VLM-powered Intelligent Mobile Communication Systems \hfill
  \normalfont\textit{\hyperlink{pub4}{\textcolor{linkblue}{[ACM UbiComp 2027]}}} \textit{\href{https://chutoutian.github.io/aquavlm.github.io/}{\textcolor{linkblue}{[Website]}}}}
\begin{itemize}
    \item Built an end-to-end interactive VR application that lets divers auto-generate and exchange context-aware messages underwater, using edge-side VLM inference over real-time multimodal capture (first-person imagery + sensor readings).
  \item Constructed an 85K-sample multimodal dataset with conversations generated by ChatGPT-4o via few-shot prompting.
  \item SFT MobileVLM-V2 (3B) on self-collected dataset and designed a hierarchical message-generation pipeline to align user intent (align rate > 80\%) while reducing inference latency (60\%-70\%).
  \item Deployed on-device model via 4-bit quantization (LLMFarm) in an iOS prototype with a customized communication protocol.
\end{itemize}


\project{AquaScope: Generative Codec-based Image Transmission on Mobile Devices \hfill
  \normalfont\textit{\hyperlink{pub2}{\textcolor{linkblue}{[ACM MobiSys 2026]}}} \textit{\href{https://chutoutian.github.io/aquascope.github.io/}{\textcolor{linkblue}{[Website]}}}}
\begin{itemize}
  \item Developed the first acoustic-based system enabling reliable image transmission between mobile devices via generative compression (TiTok, VQGAN) in resource-constrained environments.
  \item Improved bandwidth efficiency by 40\% and transmission error resilience via token clustering and fine-tuning, outperforming PNG, JPEG, and classic neural compression; reduced bit error rate (BER) from 19\% to below 2\% on average.
\end{itemize}

\par\project{GenStream: Generative Codec-based Real-time Video Streaming \hfill
  %\normalfont\textit{\textcolor{lightgray}{[Under review, ACM Multimedia 2026]}}}
  \normalfont\textit{\textcolor{lightgray}{}}}
\begin{itemize}
  \item Optimized a generative video compression model (Cosmos Tokenizer) on a single RTX 4090 GPU for real-time streaming, achieving 50\%--78\% bandwidth savings over H.264/H.265 and prior neural codecs at comparable quality.
  \item Proposed a practical token adaptation approach by exploiting video context diversity to handle network dynamics.
  \item Improved perceptual quality (LPIPS) by 38\%--59\% at matched bitrates and loss patterns compared to state-of-the-art streaming systems.
\end{itemize}


\vspace{4pt}
\roleentry{TikTok, San Jose, CA}
  {Machine Learning Engineer Intern $\cdot$ Manager: Prof. Feng Qian}
  {May 2026 -- Present }

\project{UGC Short Video Quality Assessment via MLLM\hfill
  \normalfont\textit{[Submitted to ICLR 2027]}}
\begin{itemize}
\item Built a large-scale short-video quality benchmark spanning diverse source resolutions,
      transcoding settings, and mobile devices; benchmarked SOTA MLLM scorers and systematically
      analyzed their failure modes.
\item Post-trained a context-aware MLLM scorer for short-video quality assessment via SFT with motion-artifact fusion, achieving SOTA performance on both the new benchmark and open-source dataset.
\end{itemize}

\project{Scalable Ad Video Insertion for UGC Short Videos\hfill
   \normalfont\textit{[In Preparation for ACM MMSys 2027]}}
\begin{itemize}
  \item Built an end-to-end pipeline for automatic video insertion, including video similarity matching, insertable surface detection, and composition.
\item Designed an adaptive frame sampling method for SAM3 surface detection, $8.5\times$ faster than the vanilla version, and SFT-ed Qwen3-VL embedding model for accurate brand-conflict and context similarity checks.
\end{itemize}






\iffalse
\project{360LiveCast: Low-Latency Multicast for Live 360\textdegree{} Video  \hfill
  \normalfont\textit{\hyperlink{pub5}{\textcolor{linkblue}{[IEEE MIPR 2025]}}}}
\begin{itemize}
  \item Designed a multicast framework for live 360\textdegree{} video streaming with low end-to-end latency over wireless networks.
\end{itemize}
\fi



\vspace{4pt}
\roleentry{Multimedia Lab, ByteDance Inc., San Diego, CA}
  {Research Intern $\cdot$ Manager: Dr. Shu Shi}
  {May 2024 -- Aug.\ 2025}

\project{MARS: Motion-Aware ABR for Real-Time VR \hfill
  \normalfont\textit{\hyperlink{pub1}{\textcolor{linkblue}{[Submitted to ACM MMSys 2027]}}}}
\begin{itemize}
  \item Developed a perceptually-guided bitrate adaptation algorithm for a real-time VR streaming system, achieving 15\% bandwidth reduction without degrading QoE, supported by a user study with real-world VR applications.
  \item Designed a kinematics-based state machine that synchronizes bitrate adaptation with head motion phases.
  \item Extended the system to support multiple rate control modes (CBR, VBR, CQP) and integrated motion vector for rate adaptation optimization.
\end{itemize}

\project{Trinity: Low-latency Cloud VR Gaming \hfill
  \normalfont\textit{\hyperlink{pub3}{\textcolor{linkblue}{[ACM MMSys 2026]}}} \textit{\href{https://chutoutian.github.io/trinity.github.io/}{\textcolor{linkblue}{[Website]}}}}
\begin{itemize}
  \item Validated that users exhibit different latency tolerance across action types (head, hand, body motion) using an open-source first-person shooter VR game.
  \item Developed a delayed-rendering strategy for hand motion to bridge latency gaps and a customized rendering approach to generate partial panoramas for head motion.
  \item Prototyped an end-to-end cloud VR gaming system, achieving comparable QoE to local streaming systems while reducing bandwidth cost and GPU utilization by 60\% and 40\%, respectively.
\end{itemize}



\iffalse
\vspace{4pt}
\roleentry{Institute of Wireless Communications Technology, Shanghai Jiao Tong University}
  {Graduate Research Assistant $\cdot$ Advisor: Prof.\ Ying Cui}
  {Aug.\ 2020 -- Mar.\ 2022}

\project{RL-based Adaptive Wireless Video Streaming \hfill
  \normalfont\textit{\hyperlink{pub6}{\textcolor{linkblue}{[IEEE Trans.\ Multimedia 2025]}}}}
\begin{itemize}
  \item Theoretically analyzed and optimized actor-critic reinforcement learning (RL) algorithm, accelerating training convergence; applied the RL algorithm to adaptive bitrate in video streaming.
  \item Incorporated and mathematically modeled physical-layer network information (PRB count, MCS) into the optimized RL framework, improving QoE by 14.4\%.
  \item Designed an online tuning approach via progressive neural networks to balance generalizability and environment-specific adaptation across diverse network conditions.
  \item Deployed the optimized algorithm in Tencent’s production video streaming system.
\end{itemize}
\fi

\iffalse
\project{Rate Splitting for General Multicast \hfill
  \normalfont\textit{\hyperlink{pub7}{\textcolor{linkblue}{[IEEE Trans.\ Wireless Commun.\ 2022]}}}}
\begin{itemize}
  \item Developed an optimization framework for general rate splitting in multicast systems, reducing interference and improving spectral efficiency in multi-user wireless networks.
\end{itemize}
\fi

%\project{Adaptive 360\textdegree{} Video Streaming \hfill
%  \normalfont\textit{\hyperlink{pub8}{\textcolor{linkblue}{[IEEE Trans.\ Image Process.\ 2021]}}}}
%\begin{itemize}
%  \item Designed adaptive streaming strategies for 360\textdegree{} video under perfect, imperfect, and unknown field-of-view viewing probabilities in wireless networks.
%\end{itemize}

%──────────────────────────────────────────────────────────────
\section{Selected Publications}
%──────────────────────────────────────────────────────────────

\pub{1}{\textbf{Lingzhi Zhao}, Yongqiang Gui, Yanyan Suo, Sandesh Dhawaskar Sathyanarayana, Mingyuan Wu, Bo Chen, Ruixiao Zhang, Klara Nahrstedt, and Shu Shi. MARS: Streaming How You Move with Motion-Aware ABR for Real-Time VR.}{ACM Multimedia 2026, under review.}

\pub{2}{Beitong Tian*, \textbf{Lingzhi Zhao*}, Bo Chen, Mingyuan Wu, Haozhen Zheng, Deepak Vasisht, Francis Y.\ Yan, and Klara Nahrstedt. AquaScope: Reliable Underwater Image Transmission on Mobile Devices. (*: equal contributions)}{ACM MobiSys 2026. (corresponding author)}

\pub{3}{\textbf{Lingzhi Zhao}, Yongqiang Gui, Yanyan Suo, Sandesh Dhawaskar Sathyanarayana, Ruixiao Zhang, Shu Shi, and Klara Nahrstedt. Trinity: Exploiting Latency Sensitivity to Improve Quality of Experience on Cloud VR Gaming.}{ACM MMSys 2026.}

\pub{4}{Beitong Tian*, \textbf{Lingzhi Zhao*}, Bo Chen, Haozhen Zheng, Jingcheng Yang, Mingyuan Wu, Deepak Vasisht, and Klara Nahrstedt. AquaVLM: Improving Underwater Situation Awareness with Mobile Vision Language Models. (*: equal contributions)}{ACM UbiComp 2027. (corresponding author) }

\pub{5}{\textbf{Lingzhi Zhao}, Qian Zhou, Bo Chen, and Klara Nahrstedt. 360LiveCast: A Low-Latency and Bandwidth-efficient Multicast Framework for Live 360 Video.}{IEEE MIPR 2025.}

\pub{6}{\textbf{Lingzhi Zhao}, Ying Cui, Yuhang Jia, Yunfei Zhang, and Klara Nahrstedt. Enhancing Neural Adaptive Wireless Video Streaming via Lower-Layer Information Exposure and Online Tuning.}{IEEE Transactions on Multimedia 2025.}

\pub{7}{\textbf{Lingzhi Zhao}, Ying Cui, Sheng Yang, and Shlomo Shamai (Shitz). An Optimization Framework for General Rate Splitting for General Multicast.}{IEEE Transactions on Wireless Communications 2022.}

\pub{8}{\textbf{Lingzhi Zhao}, Ying Cui, Zhi Liu, Yunfei Zhang, and Sheng Yang. Adaptive Streaming of 360 Videos with Perfect, Imperfect, and Unknown FoV Viewing Probabilities in Wireless Networks.}{IEEE Transactions on Image Processing 2021.}

%──────────────────────────────────────────────────────────────
\section{Thesis Committee}
%──────────────────────────────────────────────────────────────

\begin{tabular}{@{}l@{\hspace{16pt}}l@{}}
  \textbf{\color{darkgray}Klara Nahrstedt} \textit{\color{midgray}(Chair)} & University of Illinois Urbana-Champaign \\[2pt]
  \textbf{\color{darkgray}Francis Y.\ Yan}                                 & University of Illinois Urbana-Champaign \\[2pt]
  \textbf{\color{darkgray}Philip Brighten Godfrey}                         & University of Illinois Urbana-Champaign \\[2pt]
  \textbf{\color{darkgray}Shu Shi}                                         & ByteDance \\
\end{tabular}

\end{document}
